\section{Open problems and physical scope}\label{sec:scope}
Theorem~\ref{thm:register-headline} realizes third-level bulk exponents
already in two dimensions on factorial tori, with the gap above the
entire ground space. Theorem~\ref{thm:bulk-three} supplies a second route. Thus bulk
attainment beyond $\mathsf S_2\cup\mathsf P_2$ in this convention is
settled. What remains open is attainment under stronger definitions:
local nondegeneracy, an extremum over all KMS ground states, a unique
ground state, all-integer-size convergence, or robustness. The internal
frame still anchors the classical computation; removing such anchors is
a separate question. Section~\ref{subsec:bulk-conventions} gives the
definition separation, and Proposition~\ref{prop:critical-ceiling} gives
the distinct one-sided ceiling when a right-c.e. gap representation
has actually been established. No minimum-dimension conclusion follows.

\begin{conjecture}[Stable exponents]
A gap exponent constant on a norm-open set of fixed-local-dimension
translation-invariant affine families, with a continuously varying
critical point and a fixed gap convention, is computable.
\end{conjecture}
The explicit instability above is evidence for a distinction between
these programmed examples and robust exponents; it does not prove
this conjecture. An effectively isolated fixed point of a finite-dimensional
computable renormalisation map on an effectively compact rational domain is
computable: eliminate closed subregions certified to contain no
zero until the remaining region has the requested diameter. With
computable derivative data its linearization eigenvalues are also
computable. Any exponent formula using these eigenvalues needs
its stated nonzero and positivity conditions. Stability has not been shown to supply such a fixed
point, isolation certificate or modulus. Uncomputably complex
renormalisation flows \cite{renormalisation} make that missing
implication substantive.

\begin{conjecture}[Anchoring and dimension]
The general need for an anchor to attain two-sided inner limits,
and the necessity of dimension at least $k+1$ for $k$ nested
box-size limits in classical-forced designs, remain open.
\end{conjecture}
The restricted disjoint-axis, deterministic finite-reader result in
Section~\ref{sec:registers} gives a dimension bound for that architecture.
A fixed-width classical strip reduces to a finite state graph.
Reachable far-end labels as a function of the remaining length form
a unary regular-language problem and are effectively ultimately
periodic. A finite-valued existing limit of these labels is
eventually constant. This restricted observation does not prove a
dimension lower bound for quantum strips or all unanchored models.

Theorem~\ref{thm:robust-headline} settles the robust third level
for the anchored causal-history class: computable bottom boundary,
program origin, infinite cloning and prescribed outer limits.
Whether an unanchored translation-invariant, Bernoulli-initialized
construction has the corresponding robust scale-law range remains open.
The central equilibrium question is whether a computable finite-range
interaction can have a unique equilibrium phase with an intrinsic
uncomputable critical exponent. Theorem~\ref{thm:unique-states}
computes its local numbers but does not settle an additional asymptotic
exponent. Fault-tolerant growing readouts and noncomputable state selection
are distinct realizations, as Section~\ref{sec:states} shows.
Other questions include a genuine uncomputable correlation-length
exponent $\nu$ and a fixed local dimension independent of the
prescribed pair. A proposed structural-length construction with
size-dependent cell costs remains incomplete. Ordinary negative
finite exponents, infinite endpoints, alternative orders of limits
and all-size torus conventions are not classified by the nonnegative
attainment theorems. Changing these definitions requires new
existence and sector arguments.

\section*{How this work was done}
The authors developed the argument with AI models from more than one
developer and take responsibility for its mathematics, attribution
and presentation. The paper specifies finite classical local rules,
deterministic two-site computation tableaux, guarded reversible
counter moves, all-sector spectral estimates and arithmetic limit
arguments. For any fixed finite program these rules compile to
fixed finite matrices. No numerical diagonalization of an arbitrary
program's full many-body Hamiltonian is supplied.

The ancillary directory contains a sequential runner, readable finite
checks, recorded outputs and SHA-256 checksums. Its checks cover
penalty localization, complete finite register product spaces,
classical flag and carry rules, word interpolation identities,
wall-layer transfer matrices, digit-mixing perturbations, finite
film transport and nested-input delimiters. The bulk broadcast checks
also cover translated equality graphs, complete source fillings, the two
prefix masks, defective rulers, tie alignment, a gapless-ring witness,
reader counts and tensor sums with degenerate kernels. The independent
component analysis retains closed free paths as well as the main
counter orbit. Floating-point eigenvalues are consistency evidence,
not interval-certified proofs. The general proofs, arbitrary slow
stabilization, noncomputable limits, and the higher-dimensional
compiler are mathematical arguments, not numerically certified
limits. The ancillary README states exactly what each check covers.

The planar-register checks additionally cover translated source phases,
mask domains, literal broadcast recurrence fillings, full tensor/direct-sum
gaps and open-chain spectra. The state checks cover finite causal cones,
majority amplification, full DLR conditionals including future children,
and the constant-normalizer restriction. Restricted-response checks cover
cycle bounds, conditional variance, finite MPS contractions, dual grids,
calibrated ruler preemption and arithmetic tracking. The two inquiry checkers verify adaptive finite inverse-CDF coupling,
random-limit versus fixed-target quantifiers, the factor-two logarithmic
estimate, factorial rounding and exact ruler ties. These are finite
consistency checks, not proofs of randomized reach, noncomputability or
the absence of a computable settling bound. Robust-law checks exercise
protected outputs, drifting coordinates, all-prefix density bounds, exact
rational-power block lengths and independent finite noise cones. These computations do not prove eternal reliability, noncomputability,
or a general robust local interaction; the imported reliability theorem and
the analytic arguments have their stated scopes.

The sequential psutil runner monitors memory across the child process tree.
Individual checkers omit platform-specific peak-memory fields when those APIs
are unavailable. The numerical checks do not depend on Windows memory APIs.
Recorded process-tree memory is sampled RSS, not an operating-system historical
peak. The finite noisy examples do not certify membership in an all-time
reliability region or a universal interpreter implementation.
